\paragraph{Uniform BF16 extension.}
The following archived comparisons use the earlier host-backed reference and
physical-cost coefficients. They are not yet verified comparisons against
FAST's modeled TPU-v3 baseline. The BF16 power and interface correction in
Appendix~\ref{app:native_feasibility_repair} requires new chip recovery;
the original numerical records are retained here with that scope.
Table~\ref{tab:uniform_bf16_cnn} applies the DeepSeek comparison convention to
fresh EfficientNet-B1--B7 and UNet captures. Each classifier result covers one complete classifier
forward at its original resolution, cold checkpoint conversion/import and final
output writeback. The reference uses the same frozen matrix mapper and scaling,
900 GB/s DRAM, 32 GB/s host service, 2,048 vector lanes and 8,192 GB/s local
service as the reported DeepSeek comparison. These latter service rates are
analytical assumptions. The software control retains the reference matrix
implementations while keeping intermediate tensors in DRAM and vector
temporaries in registers.

The custom search selects a common chip and independent per-invocation tiles
from a finite recovery family. Convolution may materialize patches or read them
directly; direct reads pay 24 scalar address/mask cycles per issued input
element. The execution serializes transfers and kernels, retains all observed
classifier allocations in DRAM, and charges FP32-checkpoint-to-BF16 conversion.
Rational replay checks every original invocation, dependency, allocation,
transfer and selected service count. This is analytical recovery, not a
full-master optimality certificate or a measurement of the custom kernels.

\begin{table}[t]
\centering\small
\caption{Archived BF16 vision predictions, including preparation and writeback.
All latencies are ms. SW holds the modeled TPU-v3 resources fixed.}
\label{tab:uniform_bf16_cnn}
\begin{tabular}{@{}lrrrr@{}}
\toprule
Workload & Reference & Reference + SW & Custom & Speedup \\
\midrule
EfficientNet-B1 & 13.657 & 3.557 & 0.826 & $16.53\times$ \\
EfficientNet-B2 & 17.124 & 4.372 & 1.011 & $16.94\times$ \\
EfficientNet-B3 & 28.752 & 7.276 & 1.576 & $18.25\times$ \\
EfficientNet-B4 & 59.880 & 15.673 & 3.128 & $19.14\times$ \\
EfficientNet-B5 & 117.697 & 29.791 & 5.946 & $19.79\times$ \\
EfficientNet-B6 & 200.960 & 53.037 & 9.859 & $20.38\times$ \\
EfficientNet-B7 & 354.137 & 97.591 & 16.909 & $20.94\times$ \\
SD1.5 UNet (50 calls) & 109,439.059 & 55,112.649 & 2,712.737 & $40.34\times$ \\
\bottomrule
\end{tabular}
\end{table}

B1--B4 select 512 $16\times16$ arrays and 114 MiB SRAM; B5--B7 select 896
arrays and 185.875 MiB SRAM. Both designs have 16,384 vector lanes,
2,608 GB/s DRAM and 32,768 GB/s local service. Their host interfaces supply
512 and 256 GB/s respectively. Area/power proxies are 411.64 mm$^2$/295.55 W
and 523.44 mm$^2$/298.47 W under the earlier coefficient set. These old power
values do not establish compliance with the repaired 300-W BF16 model.
This comparison includes the explicit interface provisioning difference.
The seven fixed-resource software controls improve predicted latency by
$3.63$--$3.95\times$; full co-design gives $16.53$--$20.94\times$.
Classifiers and denoising UNets have no autoregressive prefill/decode split.

\paragraph{Same-chip feature controls.}
Table~\ref{tab:uniform_bf16_cnn_controls} removes one implementation freedom
at a time while retaining the recovered chip, all original operations and
their order, DRAM residence, cold preparation and output writeback. Direct
affine convolution reads avoid materializing patches, after paying their
address generation and masked reads. Per-invocation tiling chooses different
matrix tiles for different operations. Its control uses the best single
tile from the same five-tile menu for the whole network, while preserving
independent convolution dataflow choices. All resulting schedules pass rational
replay. These effects interact and are not additive; the controls optimize
their declared serial implementation menus, not every feature-disabled schedule.

\begin{table}[t]
\centering\small
\caption{Latency relative to the full BF16 design on the same chip.
Larger values show the cost of removing a feature.}
\label{tab:uniform_bf16_cnn_controls}
\begin{tabular}{@{}lrr@{}}
\toprule
Workload & Materialized patches only & Best single matrix tile \\
\midrule
EfficientNet-B1 & $2.00\times$ & $1.040\times$ \\
EfficientNet-B2 & $2.07\times$ & $1.049\times$ \\
EfficientNet-B3 & $2.15\times$ & $1.064\times$ \\
EfficientNet-B4 & $2.29\times$ & $1.065\times$ \\
EfficientNet-B5 & $2.38\times$ & $1.090\times$ \\
EfficientNet-B6 & $2.44\times$ & $1.096\times$ \\
EfficientNet-B7 & $2.48\times$ & $1.103\times$ \\
\bottomrule
\end{tabular}
\end{table}

\paragraph{Denoising UNet.}
The UNet result covers all 50 captured denoising calls; conditioning, the
scheduler and image decoding remain outside that declared boundary. Its chip
has 1,024 arrays, 209.75 MiB SRAM and a 128-GB/s host interface, with the same
vector, DRAM and local service rates. A paid 128-KiB GELU table is checked over
every finite BF16 input and all original tensor layouts; positional tables are
recomputed from their captured constant programs. Preparation contributes
32.508 ms and denoising with boundary transfers 2,680.229 ms. The result is an
analytical UNet design, not a complete text-to-image timing measurement.
On this same chip, restricting GELU to one banked lookup-table copy increases
latency from 2,712.737 to 7,304.729 ms ($2.693\times$); the full design chooses
the number of table replicas per operation and pays their bank capacity and
reload traffic. Materializing every convolution patch increases latency to
2,987.345 ms ($1.101\times$). In contrast, the best common matrix tile gives
2,713.233 ms, only $1.00018\times$ slower. Thus adaptive tiling contributes
little for this UNet, whereas lookup replication is consequential. Each control
retains the original 50 calls, numerical table, chip and boundary transfers;
their effects are not additive.

\paragraph{Common comparison across the suite.}
Table~\ref{tab:uniform_suite_overview} preserves the archived designs under
the BF16/FP32 arithmetic and modeled reference convention used for DeepSeek.
Each ratio compares the same complete declared task on the two chips;
task sizes differ across models and are specified below. Cold preparation
is included. The historical host-backed reference is not yet verified against
FAST's modeled TPU-v3 baseline, and these chips require re-screening under the
consistent BF16 physical-cost contract (Appendix~\ref{app:native_feasibility_repair}).
The hardware and software controls following the table identify
which changes contribute, without adding interacting speedup ratios.

\begin{table}[t]
\centering\small
\caption{Archived ratios over the earlier host-backed TPU-v3-resource reference; baseline and BF16 physical-cost reconciliation remain open. Phase ratios do not add. Dashes denote inapplicable phases.}
\label{tab:uniform_suite_overview}
\begin{tabular}{@{}lrrr@{}}
\toprule
Workload & Overall & Prefill & Decode \\
\midrule
EfficientNet-B1 & $16.53\times$ & -- & -- \\
EfficientNet-B2 & $16.94\times$ & -- & -- \\
EfficientNet-B3 & $18.25\times$ & -- & -- \\
EfficientNet-B4 & $19.14\times$ & -- & -- \\
EfficientNet-B5 & $19.79\times$ & -- & -- \\
EfficientNet-B6 & $20.38\times$ & -- & -- \\
EfficientNet-B7 & $20.94\times$ & -- & -- \\
SD1.5 UNet (50 calls) & $40.34\times$ & -- & -- \\
Llama-1 13B & $122.59\times$ & $48.00\times$ & $148.94\times$ \\
Llama-1 65B & $68.92\times$ & $25.36\times$ & $80.76\times$ \\
Llama-3 70B & $91.78\times$ & $47.14\times$ & $105.93\times$ \\
DeepSeek-V3 & $31.71\times$ & $19.85\times$ & $37.93\times$ \\
\bottomrule
\end{tabular}
\end{table}

\paragraph{Full-inference Llama designs.}
The same BF16/FP32 arithmetic and reference service rates apply to 32
independent requests per model, each producing 512 greedy output tokens.
Llama-1 retains its original 1,536-token prompt; Llama-3 uses 4,096 tokens.
Prefill emits the first output, followed by 511 decode forwards per request.
Cold checkpoint conversion/import and final backing are charged in every arm.
These are modeled extensions of the paper reference convention, not measured
TPU serving latencies. The reference remains the original host-backed
singleton schedule; the software arm holds all reference resources fixed.

Recovery selects a common chip and compatible linear batches. Between linear
stages, each member executes its original dependent operations, keeping
attention temporaries private. Gather/scatter, KV copies and all original
token dependencies are retained. Whole immutable parameters may remain in
DRAM across stages and timesteps after their first paid import. Their complete
extents are conservatively reserved in addition to the original KV/activation
peak. A bounded residency MILP selects these leases. The fixed-resource
software arm also batches linears and shares weights within each stage,
using memory-feasible complete-request cohorts and the same matrix mapper.
These finite recovery searches provide analytical executions; the separate
DeepSeek LP certificate is not transferred to them.

\begin{table}[t]
\centering\scriptsize
\caption{Uniform full-inference predictions in seconds for all 32 requests. Preparation is included in totals. SW keeps TPU-v3 resources fixed.}
\label{tab:uniform_llama_phases}
\begin{tabular}{@{}llrrrr@{}}
\toprule
Model & Phase & Reference & Reference + SW & Custom & Speedup \\
\midrule
Llama-1 13B & Prefill & 980.380 & 720.526 & 20.425 & $48.00\times$ \\
 & Decode & 17,668.410 & 1,317.840 & 118.631 & $148.94\times$ \\
 & Total & 18,663.574 & 2,053.150 & 152.245 & $122.59\times$ \\
\addlinespace
Llama-1 65B & Prefill & 4,407.777 & pending & 173.827 & $25.36\times$ \\
 & Decode & 82,145.424 & pending & 1,017.176 & $80.76\times$ \\
 & Total & 86,627.355 & pending & 1,256.908 & $68.92\times$ \\
\addlinespace
Llama-3 70B & Prefill & 15,076.868 & 10,905.168 & 319.824 & $47.14\times$ \\
 & Decode & 106,905.123 & 7,982.193 & 1,009.241 & $105.93\times$ \\
 & Total & 121,981.992 & 18,887.361 & 1,329.065 & $91.78\times$ \\
\addlinespace
\bottomrule
\end{tabular}
\end{table}

Llama-1 13B selects 512 $16\times16$ arrays, 16,384 vector lanes, 114 MiB SRAM and a 512-GB/s host interface (411.64 mm$^2$/295.55 W proxies). Its 43,350,816 logical operations complete in 1 cohort(s) of 32 requests. It retains 275 parameters occupying 13.416 GiB; the conservative DRAM bound is 63.998 GiB within 64 GiB. Prefill/decode account for 5.19\%/94.81\% of predicted time saved; preparation is separate. Fixed-resource software improves the reference by $9.09\times$; the custom design is a further $13.49\times$ faster.

Llama-1 65B selects 128 $16\times16$ arrays, 12,288 vector lanes, 74.375 MiB SRAM and a 1,024-GB/s host interface (326.38 mm$^2$/298.15 W proxies). Its 114,898,464 logical operations complete in 4 cohort(s) of 8 requests. It retains 335 parameters occupying 23.045 GiB; the conservative DRAM bound is 64.000 GiB within 64 GiB. Prefill/decode account for 4.96\%/95.03\% of predicted time saved; preparation is separate.

Llama-3 70B selects 512 $16\times16$ arrays, 16,384 vector lanes, 114 MiB SRAM and a 512-GB/s host interface (411.64 mm$^2$/295.55 W proxies). Its 133,248,544 logical operations complete in 1 cohort(s) of 32 requests. It retains 373 parameters occupying 11.241 GiB; the conservative DRAM bound is 63.997 GiB within 64 GiB. Prefill/decode account for 12.23\%/87.77\% of predicted time saved; preparation is separate. Fixed-resource software improves the reference by $6.46\times$; the custom design is a further $14.21\times$ faster.

All custom designs retain 64-GiB DRAM and use the same paid memory-interface
and local-service model as DeepSeek. These comparisons include differences
in installed interface bandwidth. Local and A100 replays check service and
logical memory; physical fragmentation and silicon timing are not measured.

\begin{table}[t]
\centering\scriptsize
\caption{Matched custom-chip controls. Extra predicted time in seconds relative to each retained design; effects are not additive.}
\label{tab:uniform_llama_controls}
\begin{tabular}{@{}llrrr@{}}
\toprule
Model & Control & Prefill extra & Decode extra & Latency ratio \\
\midrule
Llama-1 13B & No linear batching & 0.000 & 222.329 & $2.460\times$ \\
 & Singletons, no weight leases & 2.298 & 1,417.469 & $10.326\times$ \\
 & KV spill/restore & 0.115 & 137.112 & $1.901\times$ \\
 & No cross-step weight residency & 0.000 & 20.976 & $1.138\times$ \\
\addlinespace
Llama-1 65B & No linear batching & 0.000 & 1,324.921 & $2.054\times$ \\
 & Singletons, no weight leases & 6.969 & 4,909.761 & $4.912\times$ \\
 & KV spill/restore & 0.242 & 287.970 & $1.229\times$ \\
 & No cross-step weight residency & 0.139 & 94.616 & $1.075\times$ \\
\addlinespace
Llama-3 70B & No linear batching & 0.000 & 1,175.788 & $1.885\times$ \\
 & Singletons, no weight leases & 12.443 & 7,551.977 & $6.692\times$ \\
 & KV spill/restore & 0.124 & 134.051 & $1.101\times$ \\
 & No cross-step weight residency & 0.000 & 17.593 & $1.013\times$ \\
\addlinespace
\bottomrule
\end{tabular}
\end{table}


